%!TEX root =  tfg.tex
\chapter{Manual de usuario}

\begin{quotation} [Filósofo] {Karl Popper (1902--1994)}
La verdadera ignorancia no es la ausencia de conocimientos, sino el hecho de negarse a adquirirlos.
\end{quotation}

\begin{abstract}
El uso de T-Planifica es muy sencillo e intuitivo pero dejaremos claro todo lo que podemos hacer como usuario.
\end{abstract}

\section{Como crear una cuenta}

Primero se ha de crear una cuenta a través de la aplicación facilitando los datos que se muestran, nombre de usuario, contraseña, nombre, apellidos, correo electrónico, hora de inicio del horario laboral y hora de fin del horario laboral.

\begin{figure}[h]
\centering
 \includegraphics[width= \textwidth]{figures/capturas/registro.jpg} 
 \caption{Crear cuenta}
\end{figure}

\section{Como iniciar sesión}

Tras habernos registrado iremos al inicio de sesión del sistema, en el que tendremos que ingresar el nombre de usuario y la contraseña para iniciar sesión en el servicio.

\begin{figure}[h]
\centering
 \includegraphics[width= \textwidth]{figures/capturas/inicio.jpg} 
 \caption{Como iniciar sesión}
\end{figure}


\newpage

\section{Panel de control}
Tras iniciar sesión nos mostrará el panel de control, en el nos encontramos el números de usuarios en la aplicación, las notificaciones que hemos recibido, así como las invitaciones.
En la siguiente sección, podremos ver las tareas de hoy, las tareas sin empezar de hoy, así como las tareas en proceso y las tareas terminadas de hoy.
En la siguiente sección, podremos ver las tareas por equipo, seleccionando en el desplegando el equipo del que queramos ver. También podemos ver las estadísticas de mis equipos.
En la siguiente sección, podremos ver el nombre del equipo, seleccionando en el desplegando el equipo del que queramos ver, así como las estadística del equipo seleccionado en el desplegado.

\begin{figure}[h]
\centering
 \includegraphics[width= \textwidth]{figures/capturas/dashboard1.jpg} 
 \caption{Panel de control}
\end{figure}

\newpage

\begin{figure}[h]
\centering
 \includegraphics[width= \textwidth]{figures/capturas/dashboard2.jpg} 
 \caption{Panel de control}
\end{figure}

\newpage

\section{Como crear un equipo}

Como acabamos de crear una cuenta no nos encontraremos en ningún equipo para ello y visualizando el menú de la izquierda , y nos iremos a la pestaña crear equipo.

\begin{figure}[h]
\centering
 \includegraphics[width= \textwidth]{figures/capturas/crear equipo.jpg} 
 \caption{Crear un equipo}
\end{figure}

Tras presionar en crear equipo. A continuación, realizaremos el formulario dando un nombre al equipo y crearemos al equipo.

\newpage

\section{Equipos}
Para esto seleccionamos en el menú de la izquierda la pestaña Equipos, desde esta vista podremos visualizar todos los equipos creados en la aplicación. Además, si presionamos en el botón de Compañeros podremos ver los compañeros de ese equipo, así como su nombre y correo electrónico, de la misma forma lo podemos hacer con managers y asociados para ver los componentes de ambos grupos.

\begin{figure}[h]
\centering
 \includegraphics[width= \textwidth]{figures/capturas/equipos.jpg} 
 \caption{Ver equipos}
\end{figure}

\newpage

\section{Mis Equipos}
Para esto seleccionamos en el menú de la izquierda la pestaña Mis Equipos, desde esta vista podremos visualizar los equipos a los que perteneces. Además, si presionamos en el botón de Compañeros podremos ver los compañeros de ese equipo, así como su nombre y correo electrónico, de la misma forma lo podemos hacer con managers y asociados para ver los componentes de ambos grupos.

\begin{figure}[h]
\centering
 \includegraphics[width= \textwidth]{figures/capturas/misequipos.jpg} 
 \caption{Ver mis equipos}
\end{figure}

\newpage

\section{Ascender o descender a un usuario de un equipo}
Para esto seleccionamos en el menú de la izquierda la pestaña Administrados, desde esta vista podremos visualizar los equipos creados por nosotros. Además, si presionamos en el botón de Compañeros podremos ver los compañeros de ese equipo, así como su nombre y correo electrónico, de la misma forma lo podemos hacer con managers y asociados para ver los componentes de ambos grupos. Si desplegamos el selecciona una opción de ascender mánager podremos ascender un usuario de asociado a mánager. Por el contrario, si desplegamos el selecciona una opción de volver asociado podremos descender un usuario de mánager a asociado.

\begin{figure}[h]
\centering
 \includegraphics[width= \textwidth]{figures/capturas/administrar.jpg} 
 \caption{Ascender o descender a un usuario de un equipo}
\end{figure}

\newpage

\section{Como crear tareas}

Para esto seleccionamos en el menú de la izquierda la pestaña Crear Tarea, desde esta vista podremos realizar la creación de tareas.

Nos redirigirán a la siguiente vista en la que completando los siguientes campos, nombre de la tarea, descripción de la tarea, fecha inicio de la tarea , fecha para la que la tarea ha de estar terminada (deadline máximo), nivel de priorización del 1 al 4, siendo 1 la máxima prioridad y la duración de la tarea en horas. Al presionar en crear estableceremos como fecha de inicio de la tarea el día siguiente de hoy y como fecha final la establecida por el usuario.

\begin{figure}[h]
\centering
 \includegraphics[width= \textwidth]{figures/capturas/crear tarea.jpg} 
 \caption{Crear Tareas}
\end{figure}

\newpage

\section{Tareas}
Para esto seleccionamos en el menú de la izquierda la pestaña Tareas, donde podemos ver todas las tareas creadas en la aplicación. Además en esa misma vista podemos eliminar la tarea.

\begin{figure}[h]
\centering
 \includegraphics[width= \textwidth]{figures/capturas/tareas.jpg} 
 \caption{Ver tareas}
\end{figure}

\newpage

\section{Asignar una tarea de un equipo a un usuario}
Para esto seleccionamos en el menú de la izquierda la pestaña Asignar Tarea, donde podemos asignar las tareas que han sido creadas previamente. En esta vista podemos asignar la tarea al equipo que pertenezcas, así como asignar el usuario que quieres que realice la tarea. Podemos ver en la misma vista como el estado cambia de sin asignar al estado no abierta.


\begin{figure}[h]
\centering
 \includegraphics[width= \textwidth]{figures/capturas/asignacion.jpg} 
 \caption{Asignar una tarea de un equipo a un usuario}
\end{figure}

\newpage

\section{Ver perfil}
Para esto seleccionamos en la barra superior nuestro nombre de usuario. A continuación se desplegará un menú en el que podemos visualizar nuestros datos de usuario. También podremos salir de la aplicación presionando el botón de cerrar sesión.

\begin{figure}[h]
\centering
 \includegraphics[width= \textwidth]{figures/capturas/miperfil.jpg} 
 \caption{Ver perfil}
\end{figure}

\newpage

\section{Editar mi perfil}
Para hacer esto nos dirigiremos a nuestro perfil para ello volveremos a desplegar el menú Mi Perfil, y presionaremos sobre el icono de la esquina superior izquierda. Se nos redirigirá a la vista de actualizar perfil, completaremos los campos con los datos pertinentes y por último presionaremos en actualizar para reflejar estos cambios en nuestro perfil.

\begin{figure}[h]
\centering
 \includegraphics[width= \textwidth]{figures/capturas/actualizarPerfil.jpg} 
 \caption{Actualizar Perfil}
\end{figure}

\newpage

\section{Empezar tarea}

Cuando hemos asignado una tarea podemos ver seleccionar en el menú de la izquierda la pestaña Asignadas. En esta vista podremos visualizar las tareas que tenemos asignadas para realizar. Podemos seleccionar el botón de Empezar Tarea para comenzar a realizar la tarea. Además podemos ver como cambia el estado de no Abierta al estado Empezada.

\begin{figure}[h]
\centering
 \includegraphics[width= \textwidth]{figures/capturas/asignadas.jpg} 
 \caption{Empezar tarea}
\end{figure}

\newpage

\section{Terminar tarea}
Cuando hemos asignado una tarea podemos ver seleccionar en el menú de la izquierda la pestaña Asignadas. En esta vista podremos visualizar las tareas que tenemos asignadas para realizar. Podemos seleccionar el botón de Finalizar Tarea para finalizar la tarea. Además podemos ver como cambia el estado de Empezada al estado Terminada.

\begin{figure}[h]
\centering
 \includegraphics[width= \textwidth]{figures/capturas/asignadas.jpg} 
 \caption{Terminar tarea}
\end{figure}

\newpage

\section{Pausar tarea}
Cuando hemos asignado una tarea podemos ver seleccionar en el menú de la izquierda la pestaña Asignadas. En esta vista podremos visualizar las tareas que tenemos asignadas para realizar. Podemos seleccionar el botón de Pausar Tarea para pausar la tarea. Además podemos ver como cambia el estado de Empezada al estado En Pausa.

\begin{figure}[h]
\centering
 \includegraphics[width= \textwidth]{figures/capturas/asignadas.jpg} 
 \caption{Pausar tarea}
\end{figure}

\newpage


\section{Ver mis notificaciones}
Cuando se nos invite a un equipo, se nos asigne una tarea o alguien del equipo que administremos o seamos mánager se creará una notificación que se podrá ver desde la pantalla de inicio en el barra superior en el icono de una campana.

\begin{figure}[h]
\centering
 \includegraphics[width= \textwidth]{figures/capturas/notificaciones.jpg} 
 \caption{Ver mis notificaciones}
\end{figure}

\newpage

\section{Ocultar notificaciones}
Cuando se nos invite a un equipo, se nos asigne una tarea o alguien del equipo que administremos o seamos mánager se creará una notificación que se podrá ver desde la pantalla de inicio en el barra superior en el icono de una campana, si hemos leído la notificación podemos presionar el botón de ocultar notificación y no se mostrará en la vista de notificaciones.

\begin{figure}[h]
\centering
 \includegraphics[width= \textwidth]{figures/capturas/notificaciones.jpg} 
 \caption{Ver mis notificaciones}
\end{figure}

\newpage

\section{Ver mis invitaciones}
Cuando se nos invite a un equipo se creará una invitación que se podrá ver desde la pantalla de inicio en el barra superior en el icono de un lista.

\begin{figure}[h]
\centering
 \includegraphics[width= \textwidth]{figures/capturas/invitaciones.jpg} 
 \caption{Ver mis invitaciones}
\end{figure}

\newpage

\section{Crear Invitación}
Para crear una invitación a un equipo, primero desde la pantalla de inicio en el barra superior en el icono de un lista, posteriormente presionaremos sobre el icono de usuario de la esquina superior izquierda, nos llevará a una vista de crear invitación en la que podremos crear la invitación seleccionando el equipo y el correo electrónico del usuario que queremos invitar y seleccionamos el icono de la tarea con un mas.

\begin{figure}[h]
\centering
 \includegraphics[width= \textwidth]{figures/capturas/nueva invitacion.jpg} 
 \caption{Crear Invitación}
\end{figure}

\newpage

\section{Aceptar o rechazar invitación a un equipo}
Para aceptar una invitación a un equipo, primero desde la pantalla de inicio en el barra superior en el icono de un lista, posteriormente presionaremos sobre el botón ver todas las invitaciones, nos llevará a una vista de invitaciones en la que podremos aceptar o rechazar dicha invitación.

\begin{figure}[h]
\centering
 \includegraphics[width= \textwidth]{figures/capturas/aceptar invitacion.jpg} 
 \caption{Aceptar o rechazar invitación}
\end{figure}
